\part{\nibb{} Scheme}

\chapter{Why a Scheme, and which one}

\nibb{} embeds a small Scheme for \emph{orchestration only}. Scheme code
decides what to do; every effect\,---\,a file, a process, a model call\,---\,is a
Go primitive, and every primitive application crosses the guard. Scheme
never touches the network or the filesystem on its own. That gives the
harness one place to look (\adr{3}, invariant~4) and gives the user a way to
extend it without touching Go: a tool is a documented Scheme procedure
registered into the primitive table.

The language is deliberately small. It is specified in
\file{docs/scheme.md} and the interpreter implements that document and
nothing else; the conformance suite has a snippet for every form and
builtin, and a builtin that is not documented fails the suite. What follows
is a tutorial over that reference.

\begin{aside}
\textbf{On P1b.} The forms \code{letrec}, \code{case}, \code{do},
\code{when}, \code{unless}, \code{guard}, \code{with-output-to-string},
vectors and block comments arrived in Phase 1b, the first phase after the
bootstrap. Their sections are marked \textbf{[P1b]} so a reader of the
bootstrap code knows where the kernel grew.
\end{aside}

\chapter{First steps}

\section{The REPL}

\begin{sh}
$ nibb repl
nibb> (+ 1 2)
3
nibb> (define (square x) (* x x))
nibb> (map square (iota 5))
(0 1 4 9 16)
nibb> (car 1)
error: car: argument 1: expected pair, got integer 1
\end{sh}

Unbalanced input continues on the next line with a \code{...>} prompt. A
value of \emph{void}\,---\,the result of \code{define}, \code{set!},
\code{display}\,---\,is not printed. \code{nibb eval -e EXPR} evaluates one
expression; \code{nibb eval FILE} loads files. Inside a session,
\code{:eval EXPR} evaluates with \emph{user} origin, and the guard sees it
like anything else.

\section{Values}

\begin{center}
\small
\begin{tabular}{@{}llll@{}}
\toprule
Type & Written & Predicate & Notes \\
\midrule
empty list & \code{()} & \code{null?} & quote it in code: \code{'()} \\
boolean & \code{\#t} \code{\#f} & \code{boolean?} & only \code{\#f} is false \\
integer & \code{42} & \code{integer?} & int64; wraps on overflow \\
float & \code{2.5} \code{1.0} & \code{float?} & float64; always printed with a dot or exponent \\
string & \code{"hi"} & \code{string?} & immutable UTF-8; lengths and indices count runes \\
symbol & \code{foo} & \code{symbol?} & case-sensitive \\
pair / list & \code{(1 2 3)} \code{(a . b)} & \code{pair?} \code{list?} & no mutation \\
vector [P1b] & \code{\#(1 2)} & \code{vector?} & \code{vector-set!} mutates \\
procedure & \code{\#<procedure f>} & \code{procedure?} & closures and primitives alike \\
handle & \code{\#<cred:class:id>} & \code{handle?} & opaque; no builtin turns it into text \\
void & \code{\#<void>} & --- & the unspecified value \\
\bottomrule
\end{tabular}
\end{center}

\code{number?} is true for both numeric kinds. Two integers divide exactly
when the result is integral, otherwise the result is a float:

\begin{scm}
(/ 6 3)      ; 2
(/ 7 2)      ; 3.5
(/ 1.0 4)    ; 0.25
(expt 2 10)  ; 1024
(sqrt 16)    ; 4.0
(exact 2.9)  ; 2   (truncates toward zero)
(inexact 2)  ; 2.0
\end{scm}

\section{Printing}

\code{display} shows a value for humans (strings without quotes);
\code{write} shows it for machines (strings quoted and escaped so that the
reader can read them back). Both go to the interpreter's output writer,
which the host provides: stdout in the REPL, an \emph{output} event in a
session. \code{format} builds a string with \code{\textasciitilde a}
(display), \code{\textasciitilde s} (write), \code{\textasciitilde\%}
(newline) and \code{\textasciitilde\textasciitilde}:

\begin{scm}
(format "~a and ~s~%" "x" "y")   ; "x and \"y\"\n"
\end{scm}

\chapter{Definitions and procedures}

\section{define and lambda}

\begin{scm}
(define pi 3.14159)
(define (area r) (* pi r r))
(define (rest-args a . more) more)    ; dotted parameter: the rest as a list
(define variadic (lambda args args))  ; every argument as a list

(area 2)               ; 12.56636
(rest-args 1 2 3)      ; (2 3)
(variadic 1 2)         ; (1 2)
\end{scm}

A \code{define} inside a procedure body binds in that call's environment.
When \code{define} names a fresh anonymous closure, the closure takes that
name for printing and for error messages: \code{(define f (lambda (a) a))
(f)} reports \code{f: expected 1 argument(s), got 0}.

\section{Closures}

Procedures close over their defining environment. \code{set!} assigns an
\emph{existing} binding; assigning an unbound name is an error.

\begin{scm}
(define (counter)
  (let ((n 0))
    (lambda () (set! n (+ n 1)) n)))
(define c (counter))
(c) (c) (c)   ; 3
\end{scm}

\section{let, let* and named let}

\code{let} binds in parallel, \code{let*} sequentially. \emph{Named let}
is the idiomatic loop: the name is bound to the loop procedure inside the
body.

\begin{scm}
(let ((a 1) (b 2)) (+ a b))            ; 3
(let* ((a 1) (b (+ a 1))) b)           ; 2
(let loop ((i 0) (acc '()))
  (if (= i 3)
      (reverse acc)
      (loop (+ i 1) (cons i acc))))    ; (0 1 2)
\end{scm}

\section{Tail calls}

The last expression of a body, of \code{begin}, of a \code{let} body, of a
\code{cond} clause, of \code{and}/\code{or}, and either branch of \code{if}
is in \emph{tail position}. A call there does not grow the stack: a loop
written as a tail-recursive procedure runs in constant space, for a
million iterations or a billion. The test suite proves this with a
200,000-iteration loop under an 8\,MB stack cap in nine positions.

\begin{scm}
(define (count-down n)
  (if (= n 0) 'done (count-down (- n 1))))
(count-down 1000000)    ; done
\end{scm}

Mutual recursion works the same way:

\begin{scm}
(define (even2? n) (if (= n 0) #t (odd2? (- n 1))))
(define (odd2? n)  (if (= n 0) #f (even2? (- n 1))))
(even2? 1000000)        ; #t
\end{scm}

\section{letrec [P1b]}

\code{letrec} binds names that may refer to each other, for local mutual
recursion. Initialisers should be lambdas; a plain value must not depend on
a sibling that has not been evaluated yet.

\begin{scm}
(letrec ((ev? (lambda (n) (if (= n 0) #t (od? (- n 1)))))
         (od? (lambda (n) (if (= n 0) #f (ev? (- n 1))))))
  (ev? 10))   ; #t
\end{scm}

\chapter{Conditionals}

\begin{scm}
(if test then)                 ; missing else yields void
(if test then else)
(cond (test body ...)
      (test => proc)           ; pass the test's value to proc
      (else body ...))
(and a b c)                    ; short-circuit; (and) is #t
(or a b c)                     ; short-circuit; (or) is #f
\end{scm}

\code{cond} with \code{=>} is the idiom for ``look it up, and if found, use
it'':

\begin{scm}
(cond ((assq 'b '((a 1) (b 2))) => cadr)
      (else 0))    ; 2
\end{scm}

\textbf{[P1b]} \code{when} and \code{unless} evaluate a body for effect;
\code{case} dispatches on a key against lists of datums:

\begin{scm}
(when (> x 0) (display "positive") (newline))
(unless done? (display "still going"))
(case 2 ((1 2 3) 'small) ((4 5) 'medium) (else 'large))   ; small
(case 'Foo ((foo) 'lower) (else 'exact))                 ; exact -- symbols compare by name
\end{scm}

\chapter{Lists and association lists}

\section{Pairs and lists}

\begin{scm}
(cons 1 2)                 ; (1 . 2)
(list 1 2 3)               ; (1 2 3)
(car '(1 2 3))             ; 1
(cdr '(1 2 3))             ; (2 3)
(cadr '(1 2 3))            ; 2      also caar cdar cddr caddr cdddr
(length '(1 2 3))          ; 3
(append '(1) '(2 3) '())   ; (1 2 3)
(reverse '(1 2 3))         ; (3 2 1)
(list-ref '(a b c) 1)      ; b
(list-tail '(a b c) 2)     ; (c)
(iota 4)                   ; (0 1 2 3)
\end{scm}

Lists are not mutated in \nibb{} Scheme: there is no \code{set-car!}. Build
new ones.

\section{Higher-order}

\begin{scm}
(map (lambda (x) (* x x)) '(1 2 3))        ; (1 4 9)
(map + '(1 2) '(10 20 30))                 ; (11 22)   stops at the shortest
(for-each display '(1 2))                  ; prints 12, yields void
(filter odd? '(1 2 3 4))                   ; (1 3)
(fold + 0 '(1 2 3))                        ; 6         (f acc x), left to right
(fold cons '() '(1 2))                     ; ((() . 1) . 2)
(sort '(3 1 2) <)                          ; (1 2 3)   stable
(apply + 1 2 '(3 4))                       ; 10        last argument is a list
\end{scm}

\section{Membership and association lists}

\code{memq} and \code{assq} compare with \code{eq?}; \code{member} and
\code{assoc} with \code{equal?}. \code{eq?} is identity for pairs and
procedures and \emph{value} for everything else, including strings\,---\,so
\code{(eq? "a" "a")} is true, unlike in most Schemes. \code{equal?} is
structural.

\begin{scm}
(memq 'b '(a b c))                   ; (b c)
(member '(1) '((1) 2))               ; ((1) 2)
(assq 'b '((a . 1) (b . 2)))         ; (b . 2)
(assoc "k" '(("k" . 1)))             ; ("k" . 1)
\end{scm}

Association lists are \nibb's record type\,---\,JSON objects become alists,
tool metadata is an alist, \code{shell-exec} returns one. The prelude adds
\code{assoc-ref}:

\begin{scm}
(define r '(("status" . 0) ("stdout" . "ok\n") ("stderr" . "")))
(assoc-ref r "status")          ; 0
(assoc-ref r "missing")         ; #<void>
(assoc-ref r "missing" 'none)   ; none
\end{scm}

\chapter{Strings, JSON and quasiquote}

\section{Strings}

\begin{scm}
(string-length "héllo")                  ; 5   (runes)
(string-append "a" "b" "c")              ; "abc"
(substring "héllo" 1 3)                  ; "él"
(string-upcase "aé")                     ; "AÉ"
(string-trim "  x \n")                   ; "x"
(string-split "a,b,,c" ",")              ; ("a" "b" "" "c")   literal separator
(string-join '("a" "b") "-")             ; "a-b"
(string-contains? "hello" "ell")         ; #t
(string-prefix? "he" "hello")            ; #t
(string-replace "a-b-c" "-" "+")         ; "a+b+c"
(string->number "42")                    ; 42       #f if not a number
(number->string 1.5)                     ; "1.5"
(string->symbol "ab") (symbol->string 'ab)
\end{scm}

There is no character type; one-character strings serve.

\section{JSON}

\code{json-parse} and \code{json-emit} are the bridge to the model and to
tools. Objects become alists with string keys in document order, arrays
lists, \code{null} the symbol \code{null}:

\begin{scm}
(json-parse "{\"a\":1,\"b\":[true,null,2.5]}")
; (("a" . 1) ("b" #t null 2.5))

(json-emit '(("name" . "nibb") ("tags" "go" "scheme")))
; "{\"name\":\"nibb\",\"tags\":[\"go\",\"scheme\"]}"
\end{scm}

\code{json-emit} treats a proper list whose elements are all pairs with
string cars as an object, and any other proper list as an array. So
\code{'()} emits \code{[]}, and there is no way to emit \code{\{\}} in v0.
This same mapping is how a model's tool-call arguments become Scheme values.

\section{Quasiquote}

One level, with \code{,} to evaluate and \code{,@} to splice:

\begin{scm}
`(1 ,(+ 1 1) ,@(list 3 4))   ; (1 2 3 4)
\end{scm}

\chapter{Errors, output capture and vectors}

\section{Raising}

\code{(error message irritant ...)} raises an error whose text is the
message followed by the written irritants. Runtime failures use the same
kind of error with a \code{who: what} message. An error that is not caught
unwinds to the host: the REPL prints it; in a session, the model receives it
as the tool result. A primitive the guard refuses raises
\code{denied: <reason>}.

\begin{scm}
(error "boom" 1 "two")    ; error: boom 1 "two"
\end{scm}

\section{Catching [P1b]}

\code{guard} has the R7RS shape: a variable, \code{cond}-style clauses, and
a body. If the body raises, the variable is bound to the \emph{condition}\,---\,an
alist with the message and the irritants\,---\,and the clauses are tried in
order; the first true clause's value is the value of the form. No match
re-raises the original error.

\begin{scm}
(guard (e ((string-contains? (assoc-ref e 'message) "boom") 'caught)
          (else (assoc-ref e 'irritants)))
  (error "boom" 1 2))                        ; caught

(guard (e (#t (assoc-ref e 'message)))
  (car 1))                                   ; "car: argument 1: expected pair, got integer 1"

(guard (e ((assoc-ref e 'message) => string-length))
  (error "four"))                            ; 4
\end{scm}

Two things are never catchable, whatever the clauses say: a guard
\emph{denial} and a \emph{cancellation}. A tool cannot hide that it was
refused, and an abort always unwinds. The test suite checks both.

\section{Capturing output [P1b]}

\begin{scm}
(with-output-to-string (display "hi") (newline))   ; "hi\n"
\end{scm}

\section{Vectors [P1b]}

\begin{scm}
(vector 1 "a" #t)                 ; #(1 "a" #t)
(vector-ref #(10 20 30) 1)        ; 20
(let ((v (vector 1 2 3)))
  (vector-set! v 1 99)
  (vector->list v))               ; (1 99 3)
(equal? #(1 2) #(1 2))            ; #t   eq? compares identity
\end{scm}

\chapter{Primitives and the guard, from Scheme}

Every builtin, every host primitive and every registered tool is an entry in
the \emph{primitive table}. You can inspect it:

\begin{scm}
(length (primitives))                        ; ~130
(primitive-info "fs-edit")
; ((doc . "Edit a text file ...") (params . (("path" "string" "...") ...))
;  (variadic . #f) (egress . #f) (irreversible . #t) (host . #t) (origin . "builtin"))
\end{scm}

The host primitives in the bootstrap are \code{fs-read}, \code{fs-write},
\code{fs-edit}, \code{fs-list} and \code{shell-exec}. Calling one from
Scheme is calling it: the guard runs, and if it says no you get an error.

\begin{scm}
(fs-read "README.md")                   ; the contents
(fs-list ".")                           ; ("Makefile" "cmd/" ...)   dirs end in /
(fs-read ".env")                        ; error: fs-read: denied: path .env matches the credential-file denylist (.env)
(shell-exec "go" '("version"))          ; (("status" . 0) ("stdout" . "go version ...") ("stderr" . ""))
\end{scm}

\code{shell-exec} takes the program and a \emph{list} of arguments\,---\,there is
no shell, no globbing, no \code{\$VAR}. The child runs in the repository
root with an allowlisted environment; its output is capped; cancellation
kills the whole process group.

\chapter{Writing tools}

\section{register-tool}

A tool is a Scheme procedure plus metadata, registered into the primitive
table. From then on it is indistinguishable from a builtin at the dispatch
site, and the model sees it in its tool list.

\begin{scm}
(register-tool "go-test"
  (lambda (pkg) (run "go" (list "test" "-race" pkg)))
  '((doc . "Run the Go tests with the race detector for a package pattern.")
    (params . (("pkg" "string" "package pattern, e.g. ./internal/log or ./...")))))
\end{scm}

The metadata alist:

\begin{center}
\small
\begin{tabular}{@{}lp{9.5cm}@{}}
\toprule
Key & Meaning \\
\midrule
\code{doc} & one line; required; shown to the model as the tool description \\
\code{params} & list of \code{("name" "type" "doc")}; required; type is one of \code{string integer number boolean list alist vector any} \\
\code{variadic} & extra arguments allowed after \code{params}; default \code{\#f}. Prefer a list parameter: the model's schema cannot express rest arguments \\
\code{egress}, \code{irreversible} & default \code{\#f}; declare them honestly, the guard reads them \\
\bottomrule
\end{tabular}
\end{center}

Rules: the name must be new\,---\,builtins cannot be shadowed and a tool
cannot be re-registered, because the tool list must stay stable for the
whole session (\adr{9}). A registered tool's own primitive calls are guarded
in turn: a wrapper over \code{shell-exec} does not make \code{shell-exec}
less irreversible.

\section{The dev tools, read closely}

\file{dev/go.scm} is the pattern. A helper renders a \code{shell-exec}
result the way a model reads best\,---\,exit status, then the output:

\begin{scm}
(define (run cmd args)
  (let ((r (shell-exec cmd args)))
    (format "exit ~a~%~a~a"
            (assoc-ref r "status")
            (assoc-ref r "stdout")
            (assoc-ref r "stderr"))))
\end{scm}

\file{dev/repo.scm} shows a tool that needs no shell at all: find an ADR by
number by listing a directory and reading a file.

\begin{scm}
(define (adr-file n)
  (let ((prefix (format "~a-" (pad4 n))))
    (let loop ((names (fs-list "docs/adr")))
      (cond ((null? names) #f)
            ((string-prefix? prefix (car names))
             (string-append "docs/adr/" (car names)))
            (else (loop (cdr names)))))))

(register-tool "read-adr"
  (lambda (n)
    (let ((f (adr-file n)))
      (if f (fs-read f) (error "no such ADR" n))))
  '((doc . "Read architecture decision record number n from docs/adr.")
    (params . (("n" "integer" "the ADR number, e.g. 10")))))
\end{scm}

\section{init.scm and allow.toml}

\file{<repo>/.nibb/init.scm} is loaded at start, after the prelude, with
user origin and with confirms denied (nobody is listening yet), so it may
register tools and must not act. \code{load} resolves relative paths against
the repository root:

\begin{scm}
(load "dev/go.scm")
(load "dev/git.scm")
(load "dev/repo.scm")
\end{scm}

Pair each command-wrapping tool with a \emph{standing approval} in
\file{.nibb/allow.toml}, or every call prompts:

\begin{txt}
commands = ["go test", "go build", "git status", "git diff", "git grep"]
writes   = ["internal/", "docs/"]
\end{txt}

\chapter{Reference}

\section{Special forms}

\begin{center}
\small
\begin{tabular}{@{}p{7.6cm}p{6.2cm}@{}}
\toprule
Form & Meaning \\
\midrule
\code{(define name expr)} & bind in the current environment \\
\code{(define (name p ... [. rest]) body ...)} & named lambda \\
\code{(lambda (p ... [. rest]) body ...)}, \code{(lambda args body ...)} & closure \\
\code{(if test then [else])} & missing else yields void \\
\code{(let ((v e) ...) body ...)} & parallel bindings \\
\code{(let name ((v e) ...) body ...)} & named let \\
\code{(let* ((v e) ...) body ...)} & sequential bindings \\
\code{(letrec ((v e) ...) body ...)} [P1b] & mutually recursive bindings \\
\code{(begin e ...)} & sequence \\
\code{(set! name expr)} & assign an existing binding \\
\code{(quote d)}, \code{'d} & the datum \\
\code{(quasiquote d)}, \code{`d} with \code{,e} \code{,@l} & one level \\
\code{(cond (test body ...) ...} \code{(test => proc) (else ...))} & first true clause \\
\code{(case key ((d ...) body ...) ... (else ...))} [P1b] & dispatch on datums \\
\code{(and e ...)}, \code{(or e ...)} & short-circuit \\
\code{(when test e ...)}, \code{(unless test e ...)} [P1b] & conditional effect \\
\code{(do ((v init [step]) ...)} \code{(test [e ...]) cmd ...)} [P1b] & loop, constant stack \\
\code{(guard (var clause ...) body ...)} [P1b] & catch errors; denials and cancellation excepted \\
\code{(with-output-to-string e ...)} [P1b] & capture output \\
\bottomrule
\end{tabular}
\end{center}

Special-form names are reserved: binding one as a variable does not change
how the form is parsed.

\section{Builtins by group}

\begin{txt}
Numbers      + - * / = < > <= >= quotient remainder modulo abs min max
             floor ceiling round truncate exact inexact expt sqrt
             number? integer? float? zero? positive? negative? even? odd?
             number->string string->number
Equivalence  not boolean? eq? eqv? equal?
Lists        cons car cdr caar cadr cdar cddr caddr cdddr list list? pair?
             null? length append reverse list-ref list-tail iota memq member
             assq assoc map for-each filter fold sort apply
Vectors      vector vector? vector-length vector-ref vector-set!
[P1b]        vector->list list->vector
Strings      string? string-length string-append substring string=? string<?
             string-upcase string-downcase string-trim string-split string-join
             string-contains? string-prefix? string-suffix? string-replace
             string->symbol symbol->string symbol? format
Output       display write newline
Control      procedure? apply void error
JSON         json-parse json-emit
Meta         primitives primitive-info register-tool load handle?
Prelude      assoc-ref identity compose string-empty?        (Scheme)
Host         fs-read fs-write fs-edit fs-list shell-exec     (bootstrap)
\end{txt}

\section{Not in \nibb{} Scheme}

\code{set-car!}/\code{set-cdr!} and string mutation; nested quasiquote;
hygienic macros; \code{call/cc}; \code{dynamic-wind}; an \code{eval}
primitive (the kernel surface is not exposed to programs); characters;
exact rationals and bignums; source positions in runtime errors.
